% Template KLTN cho SV trường ĐHKHTN
% Liên hệ: nqminh@fit.hcmus.edu.vn
% Last update: 08/06/2016

% Chú ý: đọc các phần chú ý đóng khung của file này và chỉnh lại cho phù hợp.
% Trước khi build, xóa hết các file được tạo ra trong quá trình build trước đó, và build theo thứ tự: BIB > PDF > PDF.
% Nếu cập nhật tài liệu tham khảo, cũng cần build lại theo cách trên.

\documentclass[oneside,a4paper,14pt]{extreport}

% Font tiếng Việt
\usepackage[T5]{fontenc}
\usepackage[utf8]{inputenc}
\DeclareTextSymbolDefault{\DH}{T1}

% insert pdf
\usepackage{pdfpages}

% equation
\usepackage{breqn}
\usepackage{amsfonts}
\usepackage{bm}

% Tài liệu tham khảo
\usepackage[
	sorting=nty,
	backend=bibtex,
	defernumbers=true]{biblatex}
\usepackage[unicode]{hyperref} % Bookmark tiếng Việt
\addbibresource{References/references.bib}

\makeatletter
\def\blx@maxline{77}
\makeatother

% Chèn hình, các hình trong luận văn được để trong thư mục Images/
\usepackage{graphicx}
\usepackage{caption}
\usepackage{subcaption}
\graphicspath{ {Images/} }

% Chèn và định dạng mã nguồn
\usepackage{listings}
\usepackage{color}
\definecolor{codegreen}{rgb}{0,0.6,0}
\definecolor{codegray}{rgb}{0.5,0.5,0.5}
\definecolor{codepurple}{rgb}{0.58,0,0.82}
\definecolor{backcolour}{rgb}{0.95,0.95,0.92}
\lstdefinestyle{mystyle}{
    backgroundcolor=\color{backcolour},   
    commentstyle=\color{codegreen},
    keywordstyle=\color{magenta},
    numberstyle=\tiny\color{codegray},
    stringstyle=\color{codepurple},
    basicstyle=\footnotesize,
    breakatwhitespace=false,         
    breaklines=true,                 
    captionpos=b,                    
    keepspaces=true,                 
    numbers=left,                    
    numbersep=5pt,                  
    showspaces=false,                
    showstringspaces=false,
    showtabs=false,                  
    tabsize=2
}
\lstset{style=mystyle}

% Chèn và định dạng mã giả
\usepackage{amsmath}
\usepackage{algorithm}
\usepackage[noend]{algpseudocode}
\makeatletter
\def\BState{\State\hskip-\ALG@thistlm}
\makeatother

% chèn inline code
\usepackage{xparse}
\NewDocumentCommand{\codeword}{v}{%
    \texttt{\textcolor{blue}{#1}}%
}

% Bảng biểu
\usepackage{multirow}
\usepackage{rotating}
\usepackage{vcell}
\usepackage{array}
\usepackage{diagbox}
\usepackage{booktabs}
\usepackage{colortbl}
\newcolumntype{L}[1]{>{\raggedright\let\newline\\\arraybackslash\hspace{0pt}}m{#1}}
\newcolumntype{C}[1]{>{\centering\let\newline\\\arraybackslash\hspace{0pt}}m{#1}}
\newcolumntype{R}[1]{>{\raggedleft\let\newline\\\arraybackslash\hspace{0pt}}m{#1}}

% Đổi tên mặc định
\renewcommand{\chaptername}{Chương}
\renewcommand{\figurename}{Hình}
\renewcommand{\tablename}{Bảng}
\renewcommand{\contentsname}{Mục lục}
\renewcommand{\listfigurename}{Danh sách hình}
\renewcommand{\listtablename}{Danh sách bảng}
\renewcommand{\appendixname}{Phụ lục}

% Định dạng chapter
\usepackage{titlesec}
\titleformat{\chapter}
    [display] 
    {\normalfont\bfseries\Large}{\chaptername \ \thechapter}{10pt}{\huge}
\titlespacing*{\chapter}{0pt}{-10pt}{40pt} %khoảng cách giữa chapter và đầu trang

\titleformat{\section}
    {\normalfont\bfseries\large}{\thesection}{1em}{}

\titleformat{\subsection}
    {\normalfont\bfseries\normalsize}{\thesubsection}{1em}{}

% Dãn dòng 1.5
\usepackage{setspace}
\onehalfspacing

% Thụt vào đầu dòng
\usepackage{indentfirst}

% Canh lề
\usepackage[
    top=20mm,
    bottom=10mm,
    left=30mm,
    right=20mm,
    footskip = 15mm,
    includefoot]{geometry}

% Trang bìa
\usepackage{tikz}
\usetikzlibrary{calc}
\newcommand\HRule{\rule{\textwidth}{1pt}}

% ========================================================================================= %
% CHÚ Ý: Thông tin chung về KLTN - sinh viên điền vào đây để tự động update các trang khác  %
% ========================================================================================= %
\newcommand{\tenSV}{Nguyễn~Bảo~Long~-~Cao~Tất~Cường} % Dấu ~ là khoảng trắng không được tách (các chữ nối với nhau bằng dấu ~ sẽ nằm cùng 1 dòng
\newcommand{\mssv}{1234567}
\newcommand{\tenKL}{Meta-Learning~và~Personalization~Layer~trong~Federated~Learning} % Chú ý dấu ~ trong tên khóa luận
\newcommand{\tenGVHD}{GS.~TS.~Lê~Hoài~Bắc}
\newcommand{\tenBM}{Khoa học dữ liệu}

\begin{document}

\begin{titlepage}

    \begin{center}
    %ĐẠI HỌC QUỐC GIA THÀNH PHỐ HỒ CHÍ MINH\\
    TRƯỜNG ĐẠI HỌC KHOA HỌC TỰ NHIÊN\\
    \textbf{KHOA CÔNG NGHỆ THÔNG TIN}\\[2cm]


    { \Large \bfseries \tenSV\\[2cm] } 

    %Tên đề tài Khóa luận tốt nghiệp/Đồ án tốt nghiệp

    { \Large \bfseries Meta-Learning và Personalization Layer trong Federated Learning\\[3cm]} 


    %Chọn trong các dòng sau
    \large KHÓA LUẬN TỐT NGHIỆP CỬ NHÂN\\
    %\large ĐỒ ÁN TỐT NGHIỆP CỬ NHÂN\\
    %\large THỰC TẬP TỐT NGHIỆP CỬ NHÂN\\
    %Đưa vào dòng này nếu thuộc chương trình Chất lượng cao, hoặc lớp Cử nhân tài năng
    \large CHƯƠNG TRÌNH CHÍNH QUY\\
    %\large CHƯƠNG TRÌNH CHẤT LƯỢNG CAO\\
    %\large CHƯƠNG TRÌNH CỬ NHÂN TÀI NĂNG\\[2cm]


    \begin{tikzpicture}[remember picture, overlay]
        \draw[line width = 2pt] ($(current page.north west) + (2cm,-2cm)$) rectangle ($(current page.south east) + (-1.5cm,2cm)$);
    \end{tikzpicture}

    \vfill
    Tp. Hồ Chí Minh, tháng 02/2022

    \end{center}

    \pagebreak

    \begin{center}

    TRƯỜNG ĐẠI HỌC KHOA HỌC TỰ NHIÊN\\
    \textbf{KHOA CÔNG NGHỆ THÔNG TIN}\\[2cm]

    {\large \bfseries Nguyễn Bảo Long - 18120201\\} 
    {\large \bfseries Cao Tất Cường - 18120296\\[2cm]}

    %Tên đề tài Khóa luận tốt nghiệp/Đồ án tốt nghiệp
    { \Large \bfseries Meta-Learning và Personalization Layer trong Federated Learning\\[3cm]} 


    %Chọn trong các dòng sau
    \large KHÓA LUẬN TỐT NGHIỆP CỬ NHÂN\\
    %\large ĐỒ ÁN TỐT NGHIỆP CỬ NHÂN\\
    %Đưa vào dòng này nếu thuộc chương trình Chất lượng cao, hoặc lớp Cử nhân tài năng
    \large CHƯƠNG TRÌNH CHÍNH QUY\\[2cm]
    %\large CHƯƠNG TRÌNH CHẤT LƯỢNG CAO\\[2cm]
    %\large CHƯƠNG TRÌNH CỬ NHÂN TÀI NĂNG\\[2cm]

    \textbf{GIẢNG VIÊN HƯỚNG DẪN}\\
    \tenGVHD

    \begin{tikzpicture}[remember picture, overlay]
        \draw[line width = 2pt] ($(current page.north west) + (2cm,-2cm)$) rectangle ($(current page.south east) + (-1.5cm,2cm)$);
    \end{tikzpicture}

    \vfill
    Tp. Hồ Chí Minh, tháng 02/2022

    \end{center}

\end{titlepage}
% Sasu trang Title, các bạn chèn nhận xét gủa GVHD và GVPB. Nhận xét sẽ được giáo vụ phát sau buổi bảo vệ để các bạn đóng quyển.

\pagenumbering{roman} % Đánh số i, ii, iii, ...

%\addcontentsline{toc}{chapter}{Lời cam đoan}
%\chapter*{Lời cam đoan}
\label{reassurances}

Tôi xin cam đoan đây là công trình nghiên cứu của riêng tôi. Các số liệu và kết quả nghiên cứu trong luận văn này là trung thực và không trùng lặp với các đề tài khác.

\addcontentsline{toc}{chapter}{Lời cảm ơn}
\chapter*{Lời cảm ơn}
\label{thanks}

Tôi xin chân thành cảm ơn ...

\addcontentsline{toc}{chapter}{Đề cương chi tiết}
\chapter*{Đề cương chi tiết}
\label{proposal}

\section*{Thông tin chung}

\begin{itemize}
    \item Tên đề tài: \tenKL
    \item Giảng viên hướng dẫn: \tenGVHD
    \item Nhóm sinh viên thực hiện:
    \begin{itemize}
        \item Nguyễn Bảo Long - MSSV: 18120201
        \item Cao Tất Cường - MSSV: 18120296
    \end{itemize}

    \item Thời gian thực hiện: Từ 09/2021 đến 03/2022
    \item Loại đề tài: Nghiên cứu
\end{itemize}

\section*{Nội dung thực hiện}

\subsection*{Giới thiệu đề tài}

Trong bối cảnh bùng nổ thông tin cũng như việc đề cao tính riêng tư dữ liệu người dùng như hiện nay, các mô hình huấn luyện tập trung truyền thống dần bộc lộ nhiều điểm yếu khiến chúng không còn phù hợp. Ba điểm yếu làm cho cách tiếp cận cũ này trở nên tốn kém, không năng suất và ảnh hưởng đến quyền riêng tư của người dùng có thể kể đến:

\begin{itemize}
    \item Việc truyền dữ liệu từ máy người dùng về máy chủ để tiến hành huấn luyện tiềm ẩn nguy cơ lộ dữ liệu quan trọng của người dùng.
    \item Chi phí truyền dữ liệu từ người dùng về máy chủ để huấn luyện ngày càng lớn do lượng dữ liệu sinh ra tại thiết bị cuối ngày càng tăng cao.
    \item Cần một máy chủ thật mạnh mẽ để huấn luyện mô hình với lượng dữ liệu lớn như vậy.
\end{itemize}

Khái niệm \textit{federated learning} (FL) được Google lần đầu giới thiệu trong nghiên cứu \cite{mcmahan2017communication} với ý tưởng chính là huấn luyện mô hình máy học trên các tập dữ liệu riêng biệt được phân bố trên các thiết bị biên (được gọi là huấn luyện phân tán). Với ý tưởng này, việc triển khai mô hình máy học đến người dùng không còn gặp phải vấn đề về chi phí truyền tin, giúp bảo vệ quyền riêng tư dữ liệu và không đòi hỏi một máy chủ quá mạnh để huấn luyện mô hình. Tuy nhiên, dữ liệu của người dùng trong hệ thống thường không đồng nhất và có tính cá nhân hóa rất cao (dữ liệu Non-IID). Điều này khiến cho hiệu suất của hệ thống FL suy giảm nghiêm trọng \cite{zhao2018federated}.

Một cách ngắn gọn, hiệu suất của mô hình bị giảm là do mô hình không thích ứng nhanh được trên tập dữ liệu của người dùng. Mặt khác, các thuật toán \textit{Meta-learning} (ML) được biết đến với khả năng thích ứng nhanh trên tập dữ liệu mới \cite{hospedales2020meta}. Điều này giải quyết chính xác vấn đề dữ liệu mà FL đang gặp phải. Song song với đó, để tăng thêm tính các nhân hóa mô hình cho từng người dùng, các nghiên cứu \parencite{arivazhagan2019federated, liang2020think} đề xuất sử dụng kỹ thuật \textit{Personalization layer} (PL), giúp tăng đáng kể cả hiệu suất lẫn trải nghiệm của người dùng trong hệ thống FL. Tuy nhiên, các phương pháp tối ưu kể trên vẫn tồn tại nhiều khuyết điểm và có khả năng bù trừ cho nhau. Do đó, việc nghiên cứu về ML, PL được tiến hành và kết hợp vào hệ thống FL để đạt được hiệu suất tốt hơn.

\subsection*{Mục tiêu đề tài}

Mục tiêu chính của đề tài này bao gồm: (1) - Nghiên cứu, khảo sát các thuật toán theo hướng FL, ML, PL và (2) - Kết hợp cài đặt các thuật toán trên, giúp nâng cao hiệu suất của hệ thống FL khi đối mặt với dữ liệu Non-IID.

Việc nghiên cứu, khảo sát các thuật toán nhằm đưa ra đánh giá về ưu, nhược điểm của từng thuật toán. Từ đó, biết cách kết hợp chúng để đạt hiệu suất tốt.

Việc kết hợp cài đặt nhằm mục đích chứng minh thực nghiệm tính hiệu quả của thuật toán đề xuất khi làm việc với dữ liệu Non-IID.

\subsection*{Phạm vi đề tài}

Nghiên cứu \cite{yin2021comprehensive} chỉ ra ba hướng nghiên cứu chính khi đề cập đến một hệ thống FL: (1) - Cải thiện hiệu suất của hệ thống FL, (2) - Cải thiện khả năng bảo mật của hệ thống FL, (3) - Cải thiện vấn đề về quyền riêng tư của người dùng trong hệ thống FL.

Về việc phân loại hệ thống FL, dựa trên dữ liệu đầu vào, hệ thống FL được chia thành ba loại \cite{yin2021comprehensive}: (1) - Horizontal FL, (2) - Vertical FL, (3) - Federated transfer learning.

Về việc phân loại các kịch bản Non-IID, nghiên cứu \cite{zhu2021federated} chỉ ra bốn kịch bản chính: (1) - Phân phối thuộc tính khác nhau giữa các máy khách, (2) - Phân phối nhãn khác nhau giữa các máy khách, (3) - Phân phối thời gian khác nhau giữa các máy khách, (4) - Các kịch bản khác.

Từ đó, khóa luận này có phạm vi nghiên cứu được giới hạn và phương án giải quyết được xây dựng dựa trên ba giả định sau:

\begin{itemize}
    \item Hướng nghiên cứu: Cải thiện hiệu suất của hệ thống FL.
    \item Loại hệ thống: Môi trường thí nghiệm (bao gồm các yếu tố như số lượng người dùng, dữ liệu, cấu hình, khả năng lưu trữ của thiết bị cuối,...) tuân theo đặc trưng của hệ thống Horizontal FL.
    \item Bảo mật \& quyền riêng tư: Hệ thống đã đảm bảo tính bảo mật cũng như duy trì tốt quyền riêng tư của người dùng.
    \item Kịch bản Non-IID: Phân phối nhãn dữ liệu trên các máy khách là khác nhau.
\end{itemize}

\subsection*{Cách tiếp cận}

\textbf{Phương pháp chính.}

Hệ thống FL huấn luyện trực tiếp các mô hình máy học trên các tập dữ liệu của từng người dùng, sau đó tiến hành tổng hợp tham số của mô hình trên các máy khách này để thu được một mô hình toàn cục. Do đó, kiến trúc client-server nghiễm nhiên được nghĩ đến và trở nên phổ biến. Nghiên cứu \cite{yin2021comprehensive} nêu ra các đặc điểm chính của máy chủ và máy khách trong một hệ thống FL:

\begin{itemize}
    \item Máy chủ: Điều phối các hoạt động huấn luyện mô hình và duy trì một bộ tham số toàn cục bằng cách tổng hợp các tham số mô hình do máy khách gửi về.
    \item Máy khách: Huấn luyện mô hình học theo sự chỉ đạo của máy chủ. Chúng nhận tham số toàn cục từ máy chủ, huấn luyện mô hình trên tập dữ liệu cục bộ và gửi tham số của mô hình mới về máy chủ để tổng hợp.
\end{itemize}

Phương pháp đề xuất của khóa luận nhằm tích hợp các thuật toán ML, PL vào hệ thống nêu trên. Trong đó, các thuật toán ML được sử dụng để tạo ra một khởi tạo tốt, giúp mô hình tại máy khách hội tụ nhanh chóng (chỉ sau một hoặc một vài bước huấn luyện trên một số ít dữ liệu); các thuật toán PL được thêm vào như một phương pháp giúp tăng tính cá nhân hóa của từng mô hình máy học phân bố trên máy khách.

\textbf{Dữ liệu thực nghiệm.} Khóa luận sử dụng tập dữ liệu CIFAR-10 và tập dữ liệu MNIST trong tất cả các thí nghiệm. Cả hai tập dữ liệu đều sử dụng 75\% số điểm dữ liệu để huấn luyện và 25\% số điểm dữ liệu để kiểm thử.

\textbf{Phương pháp đối sánh.} Khóa luận sử dụng thuật toán \codeword{FedAvg} làm mô hình baseline. Để so sánh công bằng, thuật toán \codeword{FedAvgMeta} cho phép mô hình toàn cục huấn luyện theo hướng \codeword{FedAvg} được phép fine-tune một hoặc một vài bước trên một phần tập dữ liệu kiểm tra trước khi bước vào kiểm thử thực sự. Các kết quả của thuật toán đề xuất sẽ được đem so sánh với kết quả của \codeword{FedAvg} và \codeword{FedAvgMeta}.

\subsection*{Kết quả đề tài}

Sau khi tiến hành nghiên cứu, nghiên cứu này kỳ vọng đạt được những kết quả sau:

\begin{itemize}
    \item Nắm được ý tưởng huấn luyện mô hình của các thuật toán theo hướng FL, ML, PL. Cài đặt mô hình baseline.
    \item Cài đặt hệ thống FL có tích hợp ML. So sánh độ chính xác thu được với mô hình baseline.
    \item Cài đặt hệ thống FL có tích hợp ML và PL. So sánh độ chính xác thu được với mô hình baseline.
\end{itemize}

\subsection*{Kế hoạch thực hiện}

Kế hoạch thực hiện khóa luận bao gồm ba giai đoạn, được trình bày trong bảng sau:

\begin{table}[H]
    \centering
    \caption{Bảng phân chia công việc}
    \label{tab:task}
    \resizebox{\linewidth}{!}{%
    \begin{tabular}{cll} 
    \toprule
    Giai đoạn                                                                                 & \multicolumn{1}{c}{Công việc}                                                               & \multicolumn{1}{c}{Người thực hiện}  \\ 
    \hline
    \multirow{4}{*}{\begin{tabular}[c]{@{}c@{}}1 \\(01/09/2021\\- 15/09/2021)\end{tabular}}   & Tìm hiểu kiến thức nền tảng về ML                                                           & Nguyễn Bảo Long                      \\
                                                                                              & Tìm hiểu kiến thức nền tảng về FL                                                           & Cao Tất Cường                        \\
                                                                                              & Tìm hiểu kiến thức nền tảng về PL                                                           & Cao Tất Cường                        \\
                                                                                              & Trao đổi 2 mảng kiến thức                                                                   & Cả hai                               \\ 
    \hline
    \multirow{4}{*}{\begin{tabular}[c]{@{}c@{}}2\\(16/09/2021 \\- 31/01/2022)\end{tabular}}   & \begin{tabular}[c]{@{}l@{}}Cài đặt các thuật toán FedAvg, \\ FedMeta(Meta-SGD)\end{tabular} & Nguyễn Bảo Long                      \\
                                                                                              & \begin{tabular}[c]{@{}l@{}}Cài đặt các thuật toán FedAvgMeta,\\FedMeta(MAML)\end{tabular}  & Cao Tất Cường                        \\
                                                                                              & \begin{tabular}[c]{@{}l@{}}Cài đặt hệ thống FL có tích hợp ML\\và PL\end{tabular}           & Nguyễn Bảo Long                      \\
                                                                                              & Phân tích và đánh giá kết quả                                                               & Cả hai                               \\ 
    \hline
    \multirow{3}{*}{\begin{tabular}[c]{@{}c@{}}3\\(01/02/2022 \\- 25/02/2022)\\\end{tabular}} & Viết luận văn                                                                               & Nguyễn Bảo Long                      \\
                                                                                              & Làm slide thuyết trình                                                                      & Cao Tất Cường                        \\
                                                                                              & Tập thuyết trình                                                                            & Cả hai                               \\
    \bottomrule
    \end{tabular}
    }
\end{table}


\addcontentsline{toc}{chapter}{Tóm tắt}
\chapter*{Tóm tắt}
\label{abstract}

Đứng trước sự bùng nổ dữ liệu tại thiết bị biên, các phương pháp máy học truyền thống (đòi hỏi việc truyền dữ liệu từ các thiết bị cuối về một máy chủ mạnh mẽ để huấn luyện) bộc lộ nhiều nhược điểm về chi phí phần cứng và vấn đề quyền riêng tư dữ liệu của người dùng. Mặt khác, khả năng lưu trữ và tính toán tại các thiết bị biên đang ngày càng được cải thiện và nâng cao cũng như phát sự triển vượt bậc của máy học trong nhiều lĩnh vực. Thực tế này thúc đẩy việc nghiên cứu một phương pháp học tối ưu hơn về chi phí và đảm bảo quyền riêng tư dữ liệu người dùng.

Khái niệm \textit{Federated Learning} (FL) cùng thuật toán \textit{Federated Averaging} (\codeword{FedAvg}) ra đời, được xem như một giải pháp thay thế, rất phù hợp với thực tế nêu trên \cite{mcmahan2017communication}. Giải pháp này không những đạt hiệu quả gần như tương đương các phương pháp học sâu đã có, giải quyết được vấn đề chi phí phần cứng, mà còn đảm bảo được quyền riêng tư dữ liệu người dùng. Tuy nhiên, đứng trước dữ liệu không đồng nhất và có tính cá nhân hóa cao trên từng người dùng (dữ liệu Non-IID), hệ thống này bị suy giảm hiệu suất nghiêm trọng \cite{zhao2018federated}.

Khóa luận này nhằm mục đích khảo sát khái niệm \textit{Federated Learning} và vấn đề tối ưu hệ thống FL trên dữ liệu Non-IID. Bằng cách kết hợp các thuật toán huấn luyện \textit{Meta-Learning} \cite{hospedales2020meta} (ML) và sử dụng kỹ thuật \textit{Personalization Layer} \cite{zhu2021federated} (PL) vào hệ thống FL, thuật toán đề xuất của khóa luận đã đạt hiệu quả cao về độ chính xác và tính cá nhân hóa trên từng người dùng khi so sánh với thuật toán \codeword{FedAvg}, thuật toán \codeword{FedPer} \cite{arivazhagan2019federated} (tối ưu hệ thống FL bằng PL) và các thuật toán \codeword{FedMeta} \cite{chen2018federated} (tối ưu hệ thống FL bằng ML).


% Mục lục, danh sách hình, danh sách bảng
\addcontentsline{toc}{chapter}{Mục lục}
\tableofcontents
\listoffigures
\listoftables

\clearpage

\pagenumbering{arabic} % Đánh số 1, 2, 3, ...

% Các chương nội dung
\chapter{Giới thiệu}
\label{Chapter1}

\section{Đặt vấn đề \& Động lực}

Hiện nay, các thiết bị biên như điện thoại, máy tính bảng, thậm chí máy giặt, máy hút bụt thông minh có thể sinh ra lượng lớn dữ liệu trong quá trình hoạt động. Lượng dữ liệu này, nếu tận dụng được, có thể mang lại sự cải thiện rất lớn về độ chính xác cho các mô hình máy học hiện tại. Ví dụ, dữ liệu thu thập được từ bàn phím điện thoại có thể phục vụ tối ưu cho các mô hình ngôn ngữ; ảnh chụp được lưu trữ trong bộ nhớ điện thoại hoàn toàn có thể được sử dụng làm dữ liệu để huấn luyện cho mô hình nhận dạng ảnh; hay lịch sử duyệt web của người dùng có thể được dùng cho bài toán đề xuất sản phẩm. Những lý do trên trở thành một động lực to lớn, thúc đẩy việc tìm ra một phương pháp giúp tận dụng nguồn dữ liệu dồi dào này.

Việc ngày càng nhiều dữ liệu được sinh ra tại các thiết bị biên khiến cho phương pháp huấn luyện mô hình theo cách tiếp cận truyền thống (được gọi là huấn luyện tập trung) bộc lộ nhiều khuyết điểm. Ba điểm yếu khiến cho cách tiếp cận này không còn mạnh mẽ có thể kể đến: (1) - Sự vi phạm về quyền riêng tư dữ liệu, (2) - Chi phí truyền tin, (3) - Chi phí phần cứng máy chủ.

\textbf{Sự vi phạm về quyền riêng tư dữ liệu.} Phương pháp truyền thống đòi hỏi phải gửi dữ liệu người dùng về một máy chủ để tiến hành huấn luyện mô hình. Các thông tin nhạy cảm của người dùng hoàn toàn có thể bị nghe lén bởi kẻ tấn công hoặc bị khai thác khi máy chủ bị nhiễm mã độc. Điều này ảnh hưởng nghiêm trọng đến quyền riêng tư dữ liệu của người dùng - một vấn đề mà hiện nay đang nhận được rất nhiều sự quan tâm từ cả người dùng lẫn chính phủ.

\textbf{Chi phí truyền tin.} Một người dùng điện thoại thông minh giờ đây có thể thực hiện giao dịch tài chính, lướt web, xem phim ngay trên thiết bị của mình khiến cho dữ liệu được sinh ra liên tục. Một máy hút bụi thông minh được trang bị các cảm biến nên dữ liệu cũng được sinh ra liên tục trong quá trình vận hành. Điều này khiến cho dữ liệu tại thiết bị biên ngày một tăng lên. Chi phí truyền tin từ các thiết bị biên đến máy chủ để huấn luyện trở nên tốn kém và có thể gây mất thông tin, ảnh hưởng đến hiệu suất học của mô hình.

\textbf{Chi phí phần cứng máy chủ.} Sau khi dữ liệu được gửi về máy chủ, cần một cấu hình máy mạnh mẽ cùng khả năng lưu trữ lớn để có thể xử lý hết lượng dữ liệu khổng lồ trên trong thời gian giới hạn.

Việc các phương pháp tiếp cận máy học theo hướng truyền thống đang dần bộc lộ các nhược điểm về chi phí vận hành và bảo trì ngày càng cao, cũng như các mối nguy hiểm tiềm tàng có thể xảy ra đối với dữ liệu của người dùng, một lần nữa thúc đẩy việc nghiên cứu về một phương pháp huấn luyện giúp làm giảm chi phí phần cứng (sử dụng cho đường truyền và máy chủ), đồng thời đảm bảo tính riêng tư dữ liệu cho người dùng.

Khái niệm \textit{federated learning} và thuật toán \codeword{FedAvg} được đưa ra vào năm 2016 bởi Google trong nghiên cứu \cite{mcmahan2017communication} nhằm mục đích huấn luyện mô hình máy học trên các tập dữ liệu riêng biệt được phân bố trên các thiết bị biên (được gọi là huấn luyện phân tán). Do đó, \textbf{một hệ thống FL không cần một máy chủ quá mạnh để vận hành (thậm chí có thể sử dụng một máy khách để vận hành \cite{yin2021comprehensive}), không đòi hỏi chi phí truyền tin quá lớn và đảm bảo được quyền riêng tư dữ liệu của người dùng vì không diễn ra bất cứ quá trình thu thập dữ liệu từ người dùng nào (điều mà mô hình huấn luyện tập trung bắt buộc phải làm)}. Dễ thấy rằng, phần lớn quá trình tính toán được phân tán đến các thiết bị biên. Tuy nhiên, khả năng lưu trữ và tính toán tại các thiết bị này ngày càng được cải thiện, khiến cho việc huấn luyện phân tán dần trở nên khả thi và đạt hiệu quả cao hơn. Bằng chứng là sự xuất hiện của hệ thống FL trong nhiều lĩnh vực như Internet vạn vật (IoT), y tế và sức khỏe, các hệ thống đề xuất, chuỗi cung ứng,... \cite{zhou2021survey}

Bên cạnh đó, nghiên cứu \cite{zhao2018federated} chỉ ra rằng, hệ thống FL hoạt động trên nền thuật toán \codeword{FedAvg} bị giảm hiệu suất nghiêm trọng khi xử lý dữ liệu Non-IID. Tuy nhiên, dữ liệu phân bố trên máy khách là không đồng nhất và có tính cá nhân hóa tất cao. Nói cách khác, các tập dữ liệu này tuân theo phân phối Non-IID. Do đó, \textbf{khóa luận này được thực hiện nhằm nghiên cứu và cải tiến hệ thống FL để thu được kết quả cao cũng như cải thiện khả năng cá nhân hóa mô hình học cho từng người dùng trên dữ liệu Non-IID}.

\section{Phạm vi đề tài}

Nghiên cứu \cite{yin2021comprehensive} chỉ ra ba hướng nghiên cứu chính khi đề cập đến một hệ thống FL: (1) - Cải thiện hiệu suất của hệ thống FL, (2) - Cải thiện khả năng bảo mật của hệ thống FL, (3) - Cải thiện vấn đề về quyền riêng tư của người dùng trong hệ thống FL.

Về việc phân loại hệ thống FL, dựa trên dữ liệu đầu vào, hệ thống FL được chia thành ba loại \cite{yin2021comprehensive}: (1) - Horizontal FL, (2) - Vertical FL, (3) - Federated transfer learning.

Về việc phân loại các kịch bản Non-IID, nghiên cứu \cite{zhu2021federated} chỉ ra bốn kịch bản chính: (1) - Phân phối thuộc tính khác nhau giữa các máy khách, (2) - Phân phối nhãn khác nhau giữa các máy khách, (3) - Phân phối thời gian khác nhau giữa các máy khách, (4) - Các kịch bản khác.

Từ đó, khóa luận này có phạm vi nghiên cứu được giới hạn và phương án giải quyết được xây dựng dựa trên ba giả định sau:

\begin{itemize}
    \item Loại hệ thống: Môi trường thí nghiệm (bao gồm các yếu tố như số lượng người dùng, dữ liệu, cấu hình, khả năng lưu trữ của thiết bị cuối,...) tuân theo đặc trưng của hệ thống Horizontal FL.
    \item Bảo mật \& quyền riêng tư: Hệ thống đã đảm bảo tính bảo mật cũng như duy trì tốt quyền riêng tư của người dùng.
    \item Kịch bản Non-IID: Phân phối nhãn dữ liệu trên các máy khách là khác nhau.
\end{itemize}

\section{Đóng góp chính}

Các đóng góp chính của khóa luận được chia thành hai loại: đóng góp về mặt lý thuyết và đóng góp về mặt thực nghiệm.

\subsection{Đóng góp lý thuyết}

\begin{itemize}
    \item Nghiên cứu hệ thống FL và thách thức về phân phối dữ liệu mà hệ thống Horizontal FL gặp phải.
    \item Khảo sát các phương pháp tối ưu hóa hệ thống Horizontal FL trên dữ liệu dạng Non-IID. Trong đó, tập trung nghiên cứu các phương pháp theo hướng Personalized Federated Averaging \parencite{fallah2020personalized, chen2018federated} và Personalization Layer \parencite{liang2020think, arivazhagan2019federated}.
    \item Phương pháp đề xuất của khóa luận đã cho thấy khả năng đạt được độ chính xác cao hơn trong quá trình kiểm thử với hai đối tượng người dùng (người dùng cục bộ và người dùng mới) so với các phương pháp trước đó (chỉ sử dụng \codeword{FedAvg}, chỉ sử dụng Personalized Federated Averaging, hoặc chỉ sử dụng Personalization Layer).
\end{itemize}

\subsection{Đóng góp thực nghiệm}

\begin{itemize}
    \item Tổ chức bộ dữ liệu MNIST và CIFAR-10 theo hai hướng IID và Non-IID để tiến hành thí nghiệm.
    \item Cài đặt thuật toán \codeword{FedAvg}, \codeword{FedAvgMeta}, các thuật toán kết hợp giữa \codeword{FedAvg} và ML (thuật toán \codeword{FedMeta(MAML)}, \codeword{FedMeta(Meta-SGD)}).
    \item Cài đặt thuật toán kết hợp giữa \codeword{FedAvg} và PL (thuật toán \codeword{FedPer}, \codeword{LG-FedAvg}).
    \item Kết hợp các thuật toán ML và PL vào hệ thống FL: các thuật toán \codeword{FedMeta-Per}.
    \item Fine-tune các siêu tham số như số lượng máy khách tham gia huấn luyện, số bước huấn luyện cục bộ, các siêu tham số học để mô hình đạt độ chính xác tốt nhất.
\end{itemize}

\section{Bố cục}

Trong luận văn này, chương \ref{Chapter2} trình bày về tổng quan lý thuyết được sử dụng trong khóa luận, các lý thuyết này làm nền tảng cho nghiên cứu và đề xuất thuật toán; chương \ref{Chapter3} đề xuất thuật toán giúp giải quyết vấn đề vừa nêu ở chương \ref{Chapter1}; chương \ref{Chapter4} trình bày về cài đặt thực nghiệm để kiểm chứng tính hiệu quả của thuật toán; chương \ref{Chapter5} đi vào phân tích kết quả đạt được; chương \ref{Chapter6} nêu kết luận và hướng phát triển tương lai của khóa luận.


\chapter{Trình bày báo cáo}
\label{Chapter2}

\section{Quy định chung}

Báo cáo phải được trình bày ngắn gọn, rõ ràng, mạch lạc, sạch sẽ, không được tẩy xóa, có đánh số trang, đánh số bảng biểu, hình vẽ, đồ thị. 

Nội dung báo cáo được phân thành các chương. Số thứ tự của các chương, mục được đánh số bằng hệ thống số Ả-rập, không dùng số La mã. Các mục và tiểu mục được đánh số bằng các nhóm hai hoặc ba chữ số, cách nhau một dấu chấm: số thứ nhất chỉ số chương, chỉ số thứ hai chỉ số mục, số thứ ba chỉ số tiểu mục.


%Báo cáo cần dùng LaTEX để viết và trình bày theo mẫu đã được cung cấp.

 Báo cáo trình bày sử dụng khổ giấy với việc canh lề như sau: Lề trên 3 cm, lề dưới 2,5 cm, lề trái 3 cm, lề phải 2 cm. Đánh số trang ở giữa bên dưới. Đánh số trang ở giữa bên dưới.

Font chữ dùng trong báo cáo (Times New Roman) với kích cỡ (size) 13-14pt, sử dụng chế độ dãn dòng (line spacing) chế độ 1.5 lines.

%Các bảng biểu trình bày theo chiều ngang khổ giấy thì đầu bảng là lề trái của trang. 


\section{Bố cục của báo cáo}

Nội dung của báo cáo tối thiểu 50 trang khổ A4 và không nên vượt quá 100 trang (không kể các trang bìa, lời cám ơn, mục lục, tài liệu tham khảo \ldots) theo trình tự như sau:

\begin{itemize}
\item MỞ ĐẦU (thường đặt tên là ``Giới thiệu''): Trình bày lý do chọn đề tài, mục đích, đối tượng và phạm vi nghiên cứu.
Mô tả bài toán mà đề tài giải quyết.
Bài toán này có gì hay?
Tại sao lại cần giải quyết bài toán này?
Bài toán này có gì khó?
Có những hướng nào để giải quyết bài toán này?
Những hướng giải quyết trước đây có những vấn đề gì chưa giải quyết được?
Các câu hỏi nghiên cứu mà đề tài trả lời hoặc những vấn đề mà đề tài sẽ giải quyết.
Các đóng góp của đề tài.

\item TỔNG QUAN (thường đặt tên là ``Các công trình liên quan''): Phân tích đánh giá các hướng nghiên cứu đã có của các tác giả trong và ngoài nước liên quan đến đề tài; nêu những vấn đề còn tồn tại (những vấn đề nào mà các công trình khác chưa giải quyết được); chỉ ra những vấn đề mà đề tài cần tập trung, nghiên cứu giải quyết.

\item NGHIÊN CỨU THỰC NGHIỆM HOẶC LÝ THUYẾT (thường đặt tên là ``Phương pháp đề xuất''): Trình bày cơ sở lý thuyết, lý luận, giả thiết khoa học và phương pháp nghiên cứu đã được sử dụng trong đề tài.

Nếu đề xuất hướng giải quyết mới, mô hình mới thì cần mô tả chi tiết cách giải quyết của mình (chi tiết tới mức người khác có thể dựa vào phần này mà cài đặt lại được đúng hoàn toàn phương pháp của mình đề ra).

\item TRÌNH BÀY, ĐÁNH GIÁ BÀN LUẬN VỀ CÁC KẾT QUẢ (thường đặt tên là ``Kết quả thí nghiệm''): Mô tả các kết quả nghiên cứu khoa học hoặc kết quả thực nghiệm.
Đối với  đề tài ứng dụng có kết quả là sản phẩm phần mềm phải có hồ sơ thiết kế, cài đặt,\ldots theo một trong các mô hình đã học (UML,\ldots).

Thông thường cần mô tả môi trường thí nghiệm trước như sử dụng dữ liệu nào, dùng độ đo nào để đánh giá, môi trường chạy thí nghiệm (cấu hình máy nếu cần phân tích thông tin về thời gian chạy thực nghiệm). Sau đó, nêu kết quả thực nghiệm, bàn luận và giải thích kết quả.

\item KẾT LUẬN VÀ HƯỚNG PHÁT TRIỂN (thường đặt tên là ``Kết luận''): Trình bày những kết quả đạt được, những đóng góp mới và những đề xuất mới, kiến nghị về những hướng nghiên cứu tiếp theo.

\item DANH MỤC TÀI LIỆU THAM KHẢO: Chỉ bao gồm các tài liệu được trích dẫn, sử dụng và đề cập tới để bàn luận trong báo cáo.
Phần này các bạn chuẩn bị 1 file BIB để lưu các tài liệu trích dẫn.
Khi các bạn trích dẫn một tài liệu nào đó, LaTeX sẽ tự động thêm vào danh mục tài liệu tham khảo giúp các bạn.
Các bạn xem hướng dẫn cách trích dẫn ở chương sau.

\item PHỤ LỤC: Phần này bao gồm nội dung cần thiết nhằm minh họa hoặc hỗ trợ cho nội dung báo cáo như số liệu, mẫu biểu, tranh ảnh,\ldots Phụ lục không được dày hơn phần chính của báo cáo.
Nếu có công trình công bố thì để vào phần phụ lục này.
\end{itemize}

\section{Bảng biểu, hình vẽ, phương trình}

%Những qui định dưới này các bạn có thể bỏ qua hoặc đọc để hiểu thêm.
%Những định dạng này LaTeX đều tự động giúp các bạn.
%Các bạn xem hướng dẫn chi tiết hơn ở chương sau.

Việc đánh số bảng biểu, hình vẽ, phương trình phải gắn với số chương; ví dụ hình 3.4 có nghĩa là hình thứ 4 trong Chương 3.
Mọi đồ thị, bảng biểu, hình vẽ lấy từ các nguồn khác phải được trích dẫn đầy đủ.

\subsection{Bảng biểu, hình vẽ}



Đầu đề của bảng biểu ghi phía trên bảng, đầu đề của hình vẽ ghi phía dưới hình.

Thông thường, những bảng ngắn và đồ thị phải đi liền với phần nội dung đề cập tới các bảng và đồ thị này ở lần thứ nhất.
Các bảng dài có thể để ở những trang riêng nhưng cũng phải tiếp theo ngay phần nội dung đề cập tới bảng này ở lần đầu tiên.
Các bảng rộng vẫn nên trình bày theo chiều đứng dài 297mm của trang giấy, chiều rộng của trang giấy có thể hơn 210mm.
Chú ý gấp trang giấy sao cho số và đầu đề của hình vẽ hoặc bảng vẫn có thể nhìn thấy ngay mà không cần mở rộng tờ giấy.
Tuy nhiên hạn chế sử dụng các bảng quá rộng này.

Đối với những trang giấy có chiều đứng hơn 297mm (bản đồ, bản vẽ,\ldots) thì có thể để trong một phong bì cứng đính bên trong bìa sau của báo cáo.

Các hình vẽ phải sạch sẽ bằng mực đen để có thể sao chụp lại; có đánh số và ghi đầy đủ đầu đề, cỡ chữ phải bằng cỡ chữ sử dụng trong báo cáo.

Khi đề cập đến các bảng biểu và hình vẽ phải nêu rõ số của hình và bảng biểu đó, ví dụ ``... được nêu trong Bảng 4.1'' hoặc ``xem Hình 3.2'' mà không được viết ``… được nêu trong bảng dưới đây'' hoặc ``trong đồ thị của X và Y sau''.

\subsection{Phương trình toán học}

Việc trình bày phương trình toán học trên một dòng đơn hoặc dòng kép tùy ý, tuy nhiên phải thống nhất trong toàn báo cáo.

Khi ký hiệu xuất hiện lần đầu tiên thì phải giải thích và đơn vị tính phải đi kèm ngay trong phương trình có ký hiệu đó.
Nếu cần thiết, danh mục của tất cả các ký hiệu, chữ viết tắt và nghĩa của chúng cần được liệt kê và để ở phần đầu của báo cáo.

Tất cả các phương trình cần được đánh số và để trong ngoặc đơn đặt bên phía lề phải.
Nếu một nhóm phương trình mang cùng một số thì những số này cũng được để trong ngoặc, hoặc mỗi phương trình trong nhóm phương trình (5.1) có thể được đánh số là (5.1.1), (5.1.2), (5.1.3).

\section{Viết tắt}

\textbf{Không lạm dụng việc viết tắt} trong báo cáo.
Chỉ viết tắt những từ, cụm từ hoặc thuật ngữ được sử dụng nhiều lần trong báo cáo.
Không viết tắt những cụm từ  dài, những mệnh đề; không viết tắt những cụm từ ít xuất hiện trong báo cáo.
Nếu cần viết tắt những từ thuật ngữ, tên các cơ quan, tổ chức,\ldots thì được viết tắt sau lần viết thứ nhất có kèm theo chữ viết tắt trong ngoặc đơn.
Nếu báo cáo có nhiều chữ viết tắt thì phải có bảng danh mục các chữ viết tắt (xếp theo thứ tự ABC) ở phần đầu báo cáo.

Nhắc lại: \textbf{không lạm dụng việc viết tắt} trong báo cáo.
Khi các bạn sử dụng từ viết tắt, người đọc sẽ phải lật lại những phần đã đọc, để tìm lại xem từ viết tắt đó nghĩa là gì.
Việc này sẽ làm chậm tốc độ đọc và sẽ khiến người đọc khó theo dõi báo cáo của bạn hơn.
Nếu có thể, hạn chế hoàn toàn việc dùng viết tắt.

\section{Tài liệu tham khảo và cách trích dẫn}

Mọi ý kiến, khái niệm có ý nghĩa, mang tính chất gợi ý không phải của riêng tác giả và mọi tham khảo khác phải được trích dẫn và chỉ ra nguồn trong danh mục Tài liệu tham khảo của báo cáo. Nguồn được trích dẫn phải được liệt kê chính xác trong danh mục Tài liệu tham khảo.

Việc trích dẫn, tham khảo chủ yếu nhằm thừa nhận nguồn của những ý tưởng có giá trị giúp người đọc theo được mạch suy nghĩ của tác giả, không làm trở ngại việc đọc.


Không trích dẫn những kiến thức phổ biến, mọi người đều biết cũng như không làm báo cáo nặng nề với những tham khảo trích dẫn.

Nếu không có điều kiện tiếp cận được một tài liệu gốc mà phải trích dẫn thông qua một tài liệu khác thì phải nêu ra trích dẫn này, đồng thời tài liệu gốc đó không được liệt kê trong danh mục tài liệu tham khảo của báo cáo.

\chapter{Hướng dẫn sử dụng template}
\label{Chapter3}

\section{Trích dẫn tài liệu}

Dùng lệnh $\textbackslash cite$ để trích dẫn một hoặc nhiều tài liệu tham khảo.
Tài liệu tham khảo có thể là trang web~\cite{Listings,HDLVThS}, bài báo khoa học~\cite{1994-Cavnar}, sách~\cite{1984-TeX-Knuth,2006-DDien,2006-NPTV}, bài tạp chí~\cite{1989-TED} hoặc các nguồn tham khảo khác. 
Lưu ý khi trích dẫn tài liệu tham khảo, cần viết câu sao cho bỏ phần trong cặp ngoặc vuông đi thì câu vẫn đầy đủ ý nghĩa.
Ví dụ, thay vì viết ``Nghiên cứu~\cite{1989-TED} chỉ ra rằng ... '' thì nên viết ``Nghiên cứu của Zhang~\cite{1989-TED} chỉ ra rằng ...''.
Một ví dụ khác, thay vì viết ``... như trong công trình nghiên cứu~\cite{1994-Cavnar}.'' thì nên viết ``... như trong công trình nghiên cứu của Cavnar và Trenkle~\cite{1994-Cavnar}.''

\section{Chèn mã nguồn}

Để chèn mã nguồn, cần dùng package listings~\cite{Listings}:

\begin{lstlisting}
\usepackage{listings}
\end{lstlisting}

Mã nguồn có thể được chèn trực tiếp như sau:

\begin{lstlisting}[language=Python]
print "Hello, World!"
\end{lstlisting}

hoặc chèn thông qua tập tin chứa mã nguồn trong thư mục $SourceCode$ như sau:

\lstinputlisting[language=C++]{SourceCode/hello.cpp}

Để chèn mã giả, cần dùng package algorithm~\cite{Algorithm}:

\begin{lstlisting}
\usepackage{algorithm}
\end{lstlisting}

Có thể chèn mã giả vào như sau:

\begin{algorithm}
\caption{My algorithm}\label{euclid}
\begin{algorithmic}[1]
\Procedure{MyProcedure}{}
\State $\textit{stringlen} \gets \text{length of }\textit{string}$
\State $i \gets \textit{patlen}$
\BState \emph{top}:
\If {$i > \textit{stringlen}$} \Return false
\EndIf
\State $j \gets \textit{patlen}$
\BState \emph{loop}:
\If {$\textit{string}(i) = \textit{path}(j)$}
\State $j \gets j-1$.
\State $i \gets i-1$.
\State \textbf{goto} \emph{loop}.
\State \textbf{close};
\EndIf
\State $i \gets i+\max(\textit{delta}_1(\textit{string}(i)),\textit{delta}_2(j))$.
\State \textbf{goto} \emph{top}.
\EndProcedure
\end{algorithmic}
\end{algorithm}

\section{Hình ảnh}

Để chèn hình ảnh, cần dùng package graphicx~\cite{Figures}:

\begin{lstlisting}
\usepackage{graphicx}
\end{lstlisting}

Hình \ref{fig:vd1}, hình \ref{fig:vd2} là một số ví dụ về chèn hình ảnh.

\begin{figure}[htp]
\centering
\includegraphics[width=6 cm]{images/logo-khtn.png}
\caption{Hình ví dụ 1}
\label{fig:vd1}
\end{figure}

\begin{figure*}[htp]
\centering
\includegraphics[width=40 mm]{images/logo-khtn.png}
\caption{Hình ví dụ 2}
\label{fig:vd2}
\end{figure*}

\section{Bảng biểu}

Để tạo bảng biểu, tham khảo thêm tại sharelatex.com~\cite{Tables}.
Bảng~\ref{tab:vd1} là một ví dụ về bảng.
Ngoài ra, có một số tool online~\footnote{Như trang http://www.tablesgenerator.com/} có thể được dùng để tạo bảng biểu một cách trực quan.

\begin{table}[ht]
\caption{Bảng ví dụ 1}

\label{tab:vd1}%
\begin{center}
\begin{tabular}{ |p{3cm}||p{3cm}|p{3cm}|p{3cm}|  }
 \hline
 \multicolumn{4}{|c|}{Country List} \\
 \hline
 Country Name     or Area Name& ISO ALPHA 2 Code &ISO ALPHA 3 Code&ISO numeric Code\\
 \hline
 Afghanistan   & AF    &AFG&   004\\
 Aland Islands&   AX  & ALA   &248\\
 Albania &AL & ALB&  008\\
 Algeria    &DZ & DZA&  012\\
 American Samoa&   AS  & ASM&016\\
 Andorra& AD  & AND   &020\\
 Angola& AO  & AGO&024\\
 \hline
\end{tabular}
\end{center}
\end{table}

\section{Công thức}

Công thức có thể chèn vào trong cùng một dòng như $ \sqrt{a^2+b^2} $ hoặc nằm trên dòng riêng như công thức~\ref{eu_eqn}.

\begin{equation} \label{eu_eqn}
x = a_0 + \frac{1}{a_1 + \frac{1}{a_2 + \frac{1}{a_3 + a_4}}}
\end{equation}

\chapter{Cài đặt thực nghiệm}
\label{Chapter4}

\section{Mô tả dữ liệu}

Tập dữ liệu CIFAR-10 \cite{cifar-10} là tập dữ liệu hình ảnh được sủ dụng phổ biến trong việc huấn luyện các mô hình máy học hay các thuật toán thị giác máy tính. Đây là một trong các tập dữ liệu được dùng nhiều nhất trong quá trình nghiên cứu máy học. Tập dữ liệu bao gồm 60,000 ảnh màu kích thước $32\times 32$ thuộc 10 phân lớp khác nhau.

Tập dữ liệu MNIST \cite{mnist} là tập dữ liệu hình ảnh được sử dụng trong việc huấn luyện các hệ thống xử lý ảnh. Tập dữ liệu này cũng được sử dụng rộng rãi trong lĩnh vực học máy. Tập dữ liệu có tổng cộng 70,000 ảnh đen trắng các chữ số viết tay được viết bởi nhiều người.

Chúng tôi sử dụng hai tập dữ liệu MNIST và CIFAR-10 để đánh giá thuật toán đề xuất và các thuật toán được khảo sát. Bằng các đặc tính của hệ thống Horizontal FL và dữ liệu Non-IID, chúng tôi cấu hình mỗi máy khách chỉ gồm 2/10 phân lớp dữ liệu, số lượng nhãn giữa các lớp và số lượng dữ liệu giữa các máy khách là không đồng đều. Thống kê dữ liệu Non-IID được trình bày trong bảng \ref{tab:stat_noniid_data}.

\begin{table}[H]
    \caption{Thống kê trên hai tập dữ liệu MNIST và CIFAR-10 (dữ liệu Non-IID)}
    \label{tab:stat_noniid_data}
    \resizebox{\textwidth}{!}{%
    \begin{tabular}{ccccccccc}
    \toprule
    \multirow{2}{*}{Dataset} & \multirow{2}{*}{\#clients} & \multirow{2}{*}{\#samples} & \multirow{2}{*}{\#classes} & \multicolumn{4}{c}{\#samples/client} & \multirow{2}{*}{\#classes/client} \\ \cline{5-8}
                             &                            &                            &                            & min    & mean    & std     & max     &                                   \\ \hline
    MNIST                    & 50                         & 69,909                     & 10                         & 135    & 1,398   & 1,424   & 5,201   & 2                                 \\
    CIFAR-10                 & 50                         & 52,497                     & 10                         & 506    & 1,049   & 250     & 1,986   & 2                                 \\
    \bottomrule
    \end{tabular}%
    }
\end{table}

Chúng tôi cũng tiến hành các thí nghiệm của mình trên kịch bản dữ liệu IID, nơi các máy khách có phân phối giống nhau và chứa đủ dữ liệu của 10 phân lớp. Chi tiết thống kê được trình bày trong bảng \ref{tab:stat_iid_data}.

\begin{table}[H]
    \caption{Thống kê trên hai tập dữ liệu MNIST và CIFAR-10 (dữ liệu IID)}
    \label{tab:stat_iid_data}
    \resizebox{\textwidth}{!}{%
    \begin{tabular}{ccccccccc}
    \toprule
    \multirow{2}{*}{Dataset} & \multirow{2}{*}{\#clients} & \multirow{2}{*}{\#samples} & \multirow{2}{*}{\#classes} & \multicolumn{4}{c}{\#samples/client} & \multirow{2}{*}{\#classes/client} \\ \cline{5-8}
                             &                            &                            &                            & min    & mean    & std     & max     &                                   \\ \hline
    MNIST                    & 50                         & 70,000                     & 10                         & 1,395    & 1,400   & 35   & 1,645   & 10                                 \\
    CIFAR-10                 & 50                         & 60,000                     & 10                         & 1,200    & 1,200   & 0     & 1,200   & 10                                 \\
    \bottomrule
    \end{tabular}%
    }
\end{table}

Dữ liệu trên mỗi máy khách được chia làm hai tập: tập huấn luyện chiếm 75\% tổng số điểm dữ liệu và tập kiểm tra chiếm 25\% tổng số điểm dữ liệu. Theo hướng ML, chúng tôi tiếp tục chia nhỏ dữ liệu trong tập huấn luyện và tập kiểm tra thành hai tập: tập support chiếm 20\% dữ liệu và tập query chiếm 80\% dữ liệu. Như vậy, thực chất mô hình được huấn luyện trên 75\% tổng số điểm dữ liệu (tập huấn luyện), tinh chỉnh trên 5\% dữ liệu (tập support của dữ liệu kiểm tra) và kiểm thử trên 20\% tổng số điểm dữ liệu (tập query của dữ liệu kiểm tra).

Trong quá trình kiểm thử, dữ liệu kiểm thử chứa trong 50 máy khách được cấu hình để tạo ra hai loại người dùng: người dùng cũ và người dùng mới. Dữ liệu của người dùng cũ được chia như đã trình bày ở trên. Đối với người dùng mới, chúng tôi tiến hành chia lại từ 25\% dữ liệu tập kiểm tra sao cho phân phối của những người dùng này khác hoàn toàn với các phân phối đã tồn tại trước đó trong hệ thống.

Tóm lại, ký hiệu $C_{train} = \{c^{train}_1,...,c^{train}_{50}\}$ là tập máy khách dùng trong huấn luyện, $C_{test} = \{c^{test}_1,...,c^{test}_{50}\}$ là tập máy khách dùng trong kiểm thử, $N$ là tổng số điểm dữ liệu, ta có số lượng dữ liệu huấn luyện và kiểm tra lần lượt là:

\begin{equation*}
    N_{train} = \sum_{i=1}^{50} \left|c^{train}_i\right| = 0.75N
\end{equation*}
\begin{equation*}
    N_{test} = \sum_{i=1}^{50} \left|c^{test}_i\right| = 0.25N
\end{equation*}

Trong cài đặt ML, ký hiệu $c^{train}_i = \{\mathcal{D}_{train(i)}^{support}, \mathcal{D}_{train(i)}^{query}\}$, $ c^{test}_i = \{\mathcal{D}_{test(i)}^{support}, \mathcal{D}_{test(i)}^{query}\}$. Ta có số lượng dữ liệu chứa trong tập support và query của tất cả các máy khách lần lượt là:

\begin{equation*}
    N_{train/test(i)}^{support} = \left|\mathcal{D}_{train/test(i)}^{support}\right| = 0.2 \left|c_i^{train/test}\right|
\end{equation*}
\begin{equation*}
    N_{train/test(i)}^{query} = \left|\mathcal{D}_{train/test(i)}^{query}\right| = 0.8 \left|c^{train/test}_i\right|
\end{equation*}

Người dùng $c_j^{test} \in C_{test}$ được gọi là người dùng cũ nếu tồn tại người dùng $c_i^{train}\in C_{train}$ sao cho $p\left((x,y)\in c_j^{test}\right) = p\left((x,y)\in c_i^{train}\right)$. Ngược lại, $c_j^{test}$ là người dùng mới nếu $p\left((x,y)\in c_j^{test}\right) \ne p\left((x,y)\in c_i^{train}\right)$ với mọi $c_i^{train}\in C_{train}$.

\section{Phương pháp đánh giá}

Sau quá trình huấn luyện mô hình toàn cục sử dụng dữ liệu trong tập $C_{train}$, chúng tôi thực hiện đánh giá mô hình này trên dữ liệu của tập $C_{test}$ bằng cách ghi nhận lại hai thông tin \cite{chen2018federated}: (1) - Độ chính xác trong tương quan với tất cả các điểm dữ liệu, (2) - Độ chính xác trong tương quan với tất cả các máy khách.

Độ chính xác trong tương quan với tất cả các điểm dữ liệu được tính toán bằng cách duyệt qua tất cả các máy khách để thống kê số lượng mẫu dữ liệu được phân lớp đúng và tổng số mẫu dữ liệu, sau đó lấy thương của hai đại lượng này. Gọi $r_i^t, n_i$ lần lượt là số lượng mẫu được phân lớp đúng tại bước huấn luyện thứ $t$, tổng số mẫu dữ liệu trên tập dữ liệu của người dùng $c_i^{test}$ và $n$ là số người dùng tham gia kiểm thử. Độ đo này tại bước huấn luyện toàn cục thứ $t$ được tính như sau:

\begin{equation}
    AccDataPoint^t = \frac{\sum_{i=1}^{n} r_i^t}{\sum_{i=1}^{n} n_i}
\end{equation}

Độ chính xác đặt trong tương quan với tất cả các máy khách được tính bằng cách lấy trung bình cộng độ chính xác trên toàn bộ máy khách tham gia kiểm thử. Với $a_i^t$ là độ chính xác của mô hình chạy trên máy khách $c_i^{test}$, tại bước huấn luyện toàn cục thứ $t$, ta tính toán thông tin về độ chính xác và độ lệch chuẩn của $n$ người dùng như sau:

\begin{equation}
    AccClient^t = \frac{1}{n} \sum_{i=1}^n a_i^t
\end{equation}
\begin{equation}
    \sigma^t = \sqrt{\frac{1}{n} \sum_{i=1}^n {(a_i^t - AccClient^t)^2}}
\end{equation}

Trong đó, $\sigma^t$ dùng để biểu thị mức độ phân tán trên độ chính xác của người dùng kiểm thử tại bước huấn luyện thứ $t$.

Các mô hình trong hệ thống của chúng tôi được đánh giá bằng độ chính xác trong tương quan với các điểm dữ liệu sau mỗi 20 bước huấn luyện toàn cục. Độ đánh giá trong tương quan với các máy khách được tính một lần duy nhất, khi mô hình toàn cục được huấn luyện xong.

\section{Mô tả thực nghiệm}

\subsection{Kiến trúc mô hình}
\label{model_schema}

Chúng tôi sử dụng hai mô hình để rút trích đặc trưng và phân lớp dữ liệu cho tập dữ liệu CIFAR-10 và MNIST.

\textbf{CIFAR-10.} Mô hình nhận các ảnh đầu vào có kích thước $(32\times32\times3)$. Hai lớp tích chập (kernel có kích thước $(5\times5)$, số chanel lần lượt là $6$ và $16$) được sử dụng để rút trích đặc trưng. Theo sau mỗi lớp tích chập là một lớp \codeword{MaxPooing} có kích thước $(2\times2)$. Phần phân lớp gồm ba lớp tuyến tính có đầu ra lần lượt là $120$, $84$ và $10$. Các hàm kích hoạt được sử dụng là \codeword{ReLU} và \codeword{Softmax}.

\textbf{MNIST.} Mô hình nhận các ảnh đầu vào đã được làm phẳng có kích thước $(1\times784)$. Sử dụng hai lớp tuyến tính có đầu ra lần lượt là 100 và 10. Các hàm kích hoạt được sử dụng là \codeword{ReLU} và \codeword{Softmax}.

\subsection{Huấn luyện tập trung}

Chúng tôi huấn luyện tập trung dựa theo định nghĩa về hệ thống FL trong nghiên cứu \cite{yin2021comprehensive}. Theo đó, cần cấu hình tập dữ liệu $\mathcal{D}_{train}$ chứa 80\% tổng dữ liệu. Kiến trúc mạng học sâu mô tả trong phần \ref{model_schema} sẽ được huấn luyện trên tập dữ liệu này. Mô hình sau khi huấn luyện được thực thi trên tập dữ liệu $\mathcal{D}_{test}$ chứa 20\% tổng dữ liệu. Các tập $\mathcal{D}_{train}, \mathcal{D}_{test}$ có dữ liệu tuân theo phân phối đều. Kết quả về độ chính xác sau khi kiểm thử gọi là kết quả huấn luyện tập trung.

\subsection{Huấn luyện phân tán}

Trước hết chúng tôi cài đặt và huấn luyện hệ thống FL bằng thuật toán \codeword{FedAvg} để có kết quả đối sánh chính. Thuật toán này được huấn luyện trên toàn bộ dữ liệu của tập $C_{train}$ và thực hiện kiểm thử trên các tập $\mathcal{D}_{test(i)}^{query}$ của từng người dùng trong tập $C_{test}$.

Khi thuật toán \codeword{FedAvg} cập nhật xong mô hình toàn cục, trong lúc kiểm thử, chúng tôi cho phép mô hình này thực hiện tinh chỉnh một hoặc một vài bước huấn luyện trên tập dữ liệu $\mathcal{D}_{test(i)}^{support}$ của những người dùng trong tập $C_{test}$. Đây chính là ý tưởng của thuật toán \codeword{FedAvgMeta}, thuật toán sinh ra nhằm so sánh công bằng với các thuật toán \codeword{FedMeta}.

Các thuật toán \codeword{FedMeta} và \codeword{FedMeta-Per} tiến hành huấn luyện như đã trình bày tại chương \ref{Chapter3}.

Để dễ dàng trong việc tinh chỉnh các mô hình học cũng như phù hợp với phần cứng thấp của máy khách trong hệ thống Horizontal FL, chúng tôi cố định một vài tham số: mỗi bước huấn luyện toàn cục sẽ sử dụng 5 máy khách, lượng dữ liệu trong một lô (batch) là 32, số bước huấn luyện cục bộ là 1 bước, số lượng bước huấn luyện cục bộ là 300 đối với tập dữ liệu MNIST và 600 đối với tập dữ liệu CIFAR-10.

Các thuật toán được huấn luyện và kiểm thử trên dữ liệu Non-IID với hai kịch bản kiểm thử: người dùng mới và người dùng cũ. Riêng thuật toán \codeword{FedAvg}, \codeword{FedAvgMeta} được chạy trên cả dữ liệu IID lẫn Non-IID để minh họa tác động của dữ liệu Non-IID đối với hệ thống.

Các thuật toán \codeword{FedMeta} được cài đặt và sử dụng trong huấn luyện mô hình để kiểm tra khả năng thích ứng nhanh trên tập dữ liệu mới của ML khi được tích hợp vào hệ thống FL. Ngoài ra, việc này còn dùng để lấy dữ liệu so sánh với thuật toán \codeword{FedMeta-Per}.

Thuật toán \codeword{FedPer} được cài đặt sử dụng trong huấn luyện mô hình để chứng minh việc các lớp phần chung của hệ thống chưa đạt hiệu quả cao khi làm việc trên tập dữ liệu của người dùng mới. Kết quả này cũng được đem so sánh với kết của của \codeword{FedMeta-Per}.

Thuật toán \codeword{FedMeta-Per} được cài đặt sử dụng trong huấn luyện mô hình để kiểm tra khả năng thích ứng nhanh trên tập dữ liệu mới của các lớp phần chung và khả năng cá nhân hóa dựa trên từng tập dữ liệu của các lớp phần riêng. Ngoài ra, cần kiểm chứng độ chính xác của thuật toán này so với các thuật toán \codeword{FedMeta}, \codeword{FedPer} và \codeword{FedAvg}.

Tất cả các thí nghiệm của chúng tôi được giả lập bằng framework Flower 0.17.0 \cite{beutel2020flower} trên một máy chủ duy nhất. Theo đó, máy chủ và các máy khách trong hệ thống đểu sử dụng chung các tài nguyên tính toán (CPU và GPU) và giao tiếp với nhau thông qua giao thức gRPC.

Tại máy chủ, một tiến trình được đặt ra để "lắng nghe" kết nối từ các máy khách tại một cổng (port) cố định. Khi các máy khách được khởi tạo xong, máy chủ tiến hành khởi tạo tham số toàn cục bằng cách chọn ngẫu nhiên bộ tham số từ một máy khách. Tại một bước huấn luyện toàn cục, máy chủ chọn ngẫu nhiên một tập con các máy khách để gửi thông tin tới máy khách thông qua giao thức gRPC. Thông tin này giúp máy khách biết mình cần thực hiện huấn luyện hay kiểm thử mô hình. Thông tin huấn luyện bao gồm tham số toàn cục, siêu tham số huấn luyện, số bước huấn luyện cục bộ, lượng dữ liệu trong một lô dữ liệu (batch). Thông tin kiểm thử bao gồm tham số toàn cục, số bước tinh chỉnh cục bộ (nếu có), siêu tham số huấn luyện (nếu có), lượng dữ liệu trong một lô dữ liệu (nếu có). Sau khi thực hiện xong yêu cầu của máy chủ, máy khách sẽ gửi thông tin về. Trong trường hợp máy chủ yêu cầu máy khách thực hiện huấn luyện, thông tin máy chủ nhận về là các tham số cục bộ và số điểm dữ liệu tham gia huấn luyện ra tham số cục bộ đó. Máy chủ tiến hành tổng hợp tham số toàn cục từ các thông tin này. Đối với yêu cầu kiểm thử mô hình, máy chủ sẽ nhận được thông tin về độ chính xác, độ lỗi, số điểm dữ liệu được kiểm thử trên các mô hình cục bộ.

Tại máy khách, sau khi khởi tạo, máy khách kết nối đến máy chủ thông qua cổng được chỉ định trước và đợi các thông tin gửi đến từ máy chủ. Dựa trên thông tin này, máy khách sẽ thực hiện huấn luyện hay kiểm thử cục bộ. Sau đó gửi các thông tin được yêu cầu về máy chủ.


\chapter{Kết quả \& Thảo luận}
\label{Chapter5}

\section{}


\chapter{Kết luận}
\label{Chapter6}

Các tập dữ liệu không đồng nhất và có tính cá nhân hóa cao được phân bố trên các thiết bị biên của người dùng cuối đòi hỏi một hệ thống FL hiện nay cần có khả năng làm việc trên dữ liệu Non-IID. Đứng trước vấn đề này, bằng các kết hợp các thuật toán ML (thuật toán \codeword{MAML}, \codeword{Meta-SGD}) và các kỹ thuật PL (thuật toán \codeword{FedPer}, \codeword{LG-FedAvg}) vào hệ thống FL, khóa luận đề xuất thuật toán \codeword{FedMeta-Per}, một giải pháp giúp làm tăng độ chính xác lẫn tính cá nhân hóa trên từng người dùng. Trong đó, các lớp phần chung của mạng mạng học sâu được huấn luyện bằng các thuật toán ML giúp hệ thống thích ứng nhanh trên dữ liệu mới, các lớp phần riêng được duy trì tại thiết bị biên giúp làm tăng tính cá nhân hóa của mô hình học trên dữ liệu cục bộ.

Bằng thực nghiệm, khóa luận đã kiểm tra được tính hiểu quả của thuật toán đề xuất trên cả hai phương diện tăng độ chính xác và tăng tính cá nhân hóa của thuật toán đề xuất trên 50 người dùng lần lượt chứa dữ liệu của hai tập dữ liệu CIFAR-10 và MNIST so với các thuật toán có sự kết hợp của FL và ML (\codeword{FedMeta(MAML)}, \codeword{FedMeta(Meta-SGD)}) và thuật toán sử dụng kỹ thuật PL (\codeword{FedPer}). Trong đó, có thể giải thích việc đạt được kết quả cao dựa vào hai yếu tố mang tính thừa kế: (1) - Khả năng thích ứng nhanh trên tập dữ liệu mới của thuật toán đề xuất thừa hưởng từ các thuật toán ML, (2) - Khả năng cá nhân hóa cao cho từng người dùng kế thừa từ các lớp cá nhân hóa của PL và việc fine-tune dữ liệu của ML.

\textbf{Định hướng phát triển.} Thuật toán mà khóa luận đề xuất không chỉ dừng lại ở việc kết hợp bốn thuật toán kể trên mà còn có khả năng phát triển thêm dựa theo hai hướng đi lớn: (1) - Sự kết hợp các thuật toán ML theo hướng tối ưu hai cấp độ vào hệ thống FL, (2) - Việc tìm kiếm và phân cụm người dùng sao cho mỗi người dùng tìm được bộ tham số huấn luyện tốt nhất. Đối với hướng đi đầu tiên, các thuật toán ML như \codeword{FO-MAML} \cite{finn2017model}, \codeword{Reptile} \cite{nichol2018first}, \codeword{iMAML} \cite{rajeswaran2019meta} hoàn toàn có thể được tích hợp vào hệ thống. Đối với hướng đi thứ hai, cần tìm ra một độ đo tốt để việc phân cụm người dùng đạt hiệu quả cao hơn trên cả kết quả phân cụm lẫn chi phí tính toán phải bỏ ra.

Hơn thế, khóa luận còn tạo động lực cho việc khảo sát ưu, nhược điểm của các phương pháp tối ưu hệ thống hiện tại và kết hợp những ưu điểm của chúng vào cùng một hệ thống để đạt được hiệu quả tốt hơn.


% In tài liệu tham khảo
\addcontentsline{toc}{chapter}{Tài liệu tham khảo}
\printbibheading[title={Tài liệu tham khảo}]
\printbibliography[heading=subbibliography]

% \printbibliography[heading=subbibliography, title={Tiếng Việt}, resetnumbers=true]

% \DeclareNameAlias{sortname}{last-first}
% \DeclareNameAlias{default}{last-first}

% \printbibliography[heading=subbibliography, title={Tiếng Anh}, notkeyword=Viet, resetnumbers=5] 
% ===================================================================== %
% CHÚ Ý: phải gán lại resetnumbers=số tài liệu tham khảo tiếng Việt + 1 %
% ===================================================================== %

% % Phần phụ lục
\appendix

\chapter{Tìm kiếm siêu tham số}
\label{Appendix1}

Số lượng siêu tham số trong một hệ thống FL là quá lớn. Do đó, sau khi giới hạn không gian tìm kiếm thì số phép thử còn lại cũng khó có thể thực hiện vét cạn. Hệ thống FL trình bày trong khóa luận cũng không nằm ngoại lệ.

Các siêu tham số của hệ thống FL của khóa luận bao gồm: số máy khách tham gia huấn luyện trong một bước huấn luyện toàn cục (\#clients/round), số bước huấn luyện cục bộ (\#epochs), số bước huấn luyện toàn cục (\#rounds), lượng dữ liệu trong một batch dữ liệu (batch_size), số lớp phần riêng đối với các thuật toán sử dụng kỹ thuật PL (\#per_layer) và các siêu tham số học sử dụng trong tối ưu mạng học sâu bằng kỹ thuật SGD.

Để phù hợp cho phần cứng của các máy khách có cấu hình yếu trong kịch bản Horizontal FL, số bước huấn luyện cục bộ và lượng dữ liệu trong một batch dữ liệu được giữ cố định lần lượt là 1 và 32. Từ việc khảo sát các thí nghiệm FL của các nghiên cứu gần đây, số máy khách tham gia huấn luyện toàn cục được được chọn lần lượt là 2, 5 và 10 máy. Trong đó, sử dụng 5 máy khách tham gia huấn luyện cùng lúc cho kết quả cao hơn một chút so với việc sử dụng 2 hay 10 máy và tiêu tốn chi phí tính toán ở một mức chấp nhận được. Do giới hạn phần cứng và theo quan sát quá trình hội tụ, số lượng bước huấn luyện toàn cục được giữ ở mức 300 cho tập MNIST và 600 cho tập CIFAR-10.

Đối với kích thước mạng học sâu được cài đặt trong khóa luận, việc duy trì một lớp phần chung và một lớp phần riêng cho mạng neural MNIST là tất nhiên. Với mạng học sâu dùng cho tập CIFAR-10, số lớp phần riêng được chọn lần lượt là 1, 2 và 3 lớp (tính từ lớp tuyến tính cuối cùng). Kết quả chạy thực nghiệm có thấy, việc sử dụng lớp tuyến tính cuối cùng làm phần riêng và các lớp học sâu còn lại làm phần chung cho kết quả tốt nhất trên tập CIFAR-10.

Ngoại trừ siêu tham số học của từng thuật toán, các siêu tham số kể trên đều được giữ cố định trong quá trình huấn luyện. Bảng \ref{tab:fixed_hyper_param} trình bày tóm tắt các giá trị siêu tham số này.

\begin{table}[H]
    \centering
    \caption{Bảng các siêu tham số cố định của hệ thống trên MNIST và CIFAR-10}
    \label{tab:fixed_hyper_param}
    \resizebox{\linewidth}{!}{%
    \begin{tabular}{l|ccccc} 
    \toprule
             & \#clients/round    & \#epochs           & \#rounds & batch\_size         & \multicolumn{1}{l}{\#per\_layer}  \\ 
    \hline
    MNIST    & \multirow{2}{*}{5} & \multirow{2}{*}{1} & 300      & \multirow{2}{*}{32} & \multirow{2}{*}{1}                \\
    CIFAR-10 &                    &                    & 600      &                     &                                   \\
    \bottomrule
    \end{tabular}
    }
\end{table}

Các siêu tham số học được tìm kiếm trong khoảng $(10^{-5}, 0.01)$ cho từng thuật toán. Kết quả tìm kiếm được trình bày trong Bảng \ref{tab:hyper_param}.

\begin{table}[H]
    \centering
    \caption{Bảng siêu tham số được sử dụng cho từng thuật toán}
    \label{tab:hyper_param}
    \resizebox{\linewidth}{!}{%
    \begin{tabular}{l|cc} 
    \toprule
    \begin{tabular}[c]{@{}l@{}}\\\end{tabular} & CIFAR-10           & MNIST                                             \\ 
    \hline
    FedAvg, FedAvgMeta                         & $10^{-4}$          & $10^{-5}$                                      \\
    FedPer, FedPerMeta                         & $10^{-4}$          & $10^{-5}$                                      \\
    FedMeta(MAML) ($\alpha,\beta$)             & $(0.01, 0.001)$    & $(0.001, 0.001)$                                \\
    FedMeta(Meta-SGD)($\alpha,\beta$)          & $(0.001, 0.001)$   & $(0.001, 5\times 10^{-4})$  \\
    FedMeta-Per(MAML)($\alpha,\beta$)          & $(0.001, 0.005)$   & $(0.001, 0.001)$                                \\
    FedMeta-Per(Meta-SGD)($\alpha,\beta$)      & $(0.01,0.01)$      & $(0.001, 5\times 10^{-4})$  \\
    \bottomrule
    \end{tabular}
    }
\end{table}
% \chapter{Ngữ pháp tiếng Nôm}
\label{Appendix2}

Đây là phụ lục 2.

\end{document} 